\chapter{Metodología}
\label{ch:metodologia}

El presente capítulo describe el enfoque metodológico utilizado para estimar 
el aporte de las cooperativas de ahorro y crédito al producto interno bruto 
de Chile, incluyendo las fuentes de información utilizadas, los métodos de 
estimación aplicados y el panel interactivo desarrollado como herramienta de 
visualización de los resultados. La sección~\ref{sec:aspectos-conceptuales} 
sintetiza la estrategia de estimación derivada del marco conceptual del 
Capítulo~\ref{ch:marco}; la sección~\ref{sec:fuentes-informacion} presenta las 
fuentes de información disponibles y su diagnóstico de cobertura; la 
sección~\ref{sec:metodos-estimacion} detalla los métodos de estimación, 
incluyendo la integración de fuentes, la operacionalización de las variables 
del SCN y el modelo de cálculo del VAB; y la sección~\ref{sec:panel} presenta 
el panel interactivo de visualización desarrollado como complemento a los 
resultados.

\section{Aspectos conceptuales}
\label{sec:aspectos-conceptuales}

\subsection{Estrategia de estimación}
\label{subsec:marco-referencia}

El presente estudio se apoya en el marco conceptual desarrollado en el 
Capítulo~\ref{ch:marco}. Como se estableció en las 
secciones~\ref{sec:cac} y~\ref{sec:cuentas-satelite}, las CAC son unidades 
institucionales del sector financiero cuya producción de mercado es 
directamente medible, y los tres marcos metodológicos vinculantes ---SCN 
2025, Manual ONU-TSE \cite{vereinte_nationen_satellite_2018} y Manual CIRIEC 
\cite{barea_tejeiro_manual_2007}--- convergen en un núcleo común de variables: 
producción (P1), consumo intermedio (P2), valor agregado bruto (B1g), 
remuneraciones (D1) y excedente de explotación (B2g).

La estrategia de estimación se deriva directamente de ese marco. Dado que 
las CAC generan producción de mercado, P1 se obtiene desde los ingresos 
operacionales de sus estados financieros, sin valoración al costo 
(véase sección~\ref{subsec:cac-scn}). Dado que el consumo intermedio no es 
observable en las fuentes administrativas chilenas ---por la exención 
tributaria del artículo~78 del DFL N\textsuperscript{o}~5/2003---, P2 se 
estima mediante un coeficiente técnico derivado de la Matriz Insumo-Producto, 
procedimiento admitido por el SCN cuando no existe declaración tributaria 
desagregada (véase sección~\ref{subsec:cac-scn}).

La metodología combina tres componentes principales:

\begin{enumerate}
    \item[a)] \textbf{Identificación y caracterización de organizaciones}: 
    extracción y validación de los RUTs de cooperativas de ahorro y crédito 
    a partir de registros administrativos disponibles.
    \item[b)] \textbf{Integración y estandarización de fuentes de datos}: 
    articulación de bases de datos provenientes de distintos organismos 
    públicos mediante el RUT como identificador común.
    \item[c)] \textbf{Estimación del aporte económico}: construcción de 
    variables compatibles con el SCN y cálculo del VAB mediante un modelo 
    de estimación indirecta, cuyos fundamentos conceptuales se establecen 
    en la sección~\ref{subsec:cac-scn}.
\end{enumerate}


\section{Fuentes de información}
\label{sec:fuentes-informacion}

La construcción de una cuenta satélite de las cooperativas de ahorro y crédito exige, como condición previa, identificar qué datos son necesarios según el marco metodológico internacional y en qué medida están disponibles en Chile. Esta sección presenta primero el inventario de fuentes administrativas públicas utilizadas en este estudio, luego sistematiza los requerimientos de información establecidos por los manuales internacionales, y finalmente contrasta ambos para identificar brechas y decisiones metodológicas resultantes.

\subsection{Inventario de fuentes administrativas chilenas}
\label{subsec:inventario-fuentes}
El análisis integra tres fuentes de datos administrativos, todas disponibles en 
formato digital y accesibles para investigación:

\begin{itemize}
    \item \textbf{División de Asociatividad y Cooperativas (DAES)}: registro oficial 
    de cooperativas activas en Chile, incluyendo tipo, región y RUT. Constituye la 
    fuente primaria para la identificación del universo de CAC no supervisadas por la 
    CMF. Cada cooperativa publica en el portal DAES sus estados financieros anuales 
    en formato PDF, de los cuales se extrajeron manualmente las variables económicas 
    necesarias para la cuenta satélite.
    \item \textbf{Comisión para el Mercado Financiero (CMF)}: provee información 
    financiera de las 7 CAC bajo su supervisión a través de la plataforma BEST y las 
    memorias anuales individuales de cada entidad. Constituye la fuente primaria para 
    el segmento CMF de la estimación.
    \item \textbf{Servicio de Impuestos Internos (SII)}: provee series de tiempo de 
    ventas por tramo y número de trabajadores por RUT para personas jurídicas. Se 
    utilizó como fuente complementaria para la caracterización institucional de las 
    CAC presentes en sus registros.
\end{itemize}

El Cuadro~\ref{tab:cobertura_fuentes} resume la cobertura de cada fuente en términos 
de entidades, período y variables disponibles.

\begin{table}[H]
\centering
\scriptsize
\setlength{\tabcolsep}{4pt}
\caption{Cobertura de las fuentes administrativas utilizadas}
\label{tab:cobertura_fuentes}
\begin{tabular*}{\textwidth}{@{\extracolsep{\fill}}
  p{2.5cm} p{2cm} p{2.5cm} p{6cm}}
\toprule
\textbf{Fuente} & \textbf{N entidades} & \textbf{Período} & \textbf{Variables principales} \\
\midrule
DAES & 35 CAC & 2014--2024 & P1, D1, activos, patrimonio, depósitos captados (F2), 
  préstamos otorgados (F4); extraídos de estados financieros en PDF \\
\addlinespace
CMF  & 7 CAC  & 2013--2025 & P1, D1, activos, patrimonio, depósitos captados (F2), 
  préstamos otorgados (F4); obtenidos de plataforma BEST y memorias anuales \\
\addlinespace
SII  & 12 CAC & Complementario & Tramo de ventas y número de trabajadores por RUT; 
  utilizado para caracterización institucional \\
\bottomrule
\end{tabular*}
\\[0.8em]
{\footnotesize \textit{Fuente: Elaboración propia.}}
\end{table}

Este inventario dialoga directamente con el resultado tecnológico comprometido en el formulario FONDEF (Anexo ~\ref{anx:resultados tecnologicos}), que declara la integración de al menos 15 fuentes de datos públicos para el conjunto de la economía social. Frente a esa declaración de amplitud, esta memoria aporta la profundidad que el formulario no desarrollaba: selecciona las fuentes pertinentes al caso piloto CAC y documenta, para cada una, su cobertura efectiva de entidades, período y variables (Cuadro \ref{sec:aspectos-conceptuales}), su método de extracción y las brechas resultantes (Sección \ref{subsec:diagnostico-disponibilidad-ch4}). Con ello se responde, a escala piloto, a una de las observaciones del panel evaluador del concurso: la ausencia de procesos explícitos de verificación y normalización de los datos integrados.

\subsection{Estadística descriptiva del panel}
\label{subsec:estadistica-descriptiva}
El panel consolidado comprende dos segmentos con características estructuralmente 
distintas. El segmento DAES cubre 35 CAC no supervisadas por la CMF para el período 
2014--2024, con un total de 129 observaciones para la variable de producción total 
(P1). El segmento CMF cubre las 7 CAC supervisadas para el período 2013--2025, con 
91 observaciones por variable. La diferencia de escala entre ambos segmentos es 
sustantiva: la producción total media del segmento CMF (MM\$~42.668) supera en más 
de 50 veces a la del segmento DAES (MM\$~828), lo que refleja la concentración del 
sector en un número reducido de cooperativas de gran tamaño.

Al interior de cada segmento, la dispersión es elevada. En el segmento DAES, la 
mediana de P1 (MM\$~392) es menos de la mitad de la media (MM\$~828), evidenciando 
la presencia de cooperativas de mayor tamaño que sesgan la distribución hacia arriba. 
Este patrón se repite en activos, patrimonio, depósitos captados (F2) y préstamos 
otorgados (F4). En el segmento CMF, la asimetría es aún más pronunciada: la mediana 
de activos (MM\$~54.404) representa apenas el 15\% de la media (MM\$~359.223), 
reflejo de la posición dominante de Coopeuch dentro del segmento supervisado. El 
Cuadro~\ref{tab:estadistica_descriptiva} presenta el resumen estadístico completo 
para ambos segmentos.

\begin{table}[H]
\centering
\scriptsize
\setlength{\tabcolsep}{4pt}
\caption{Estadística descriptiva de variables principales por segmento. 
Cifras en millones de pesos corrientes (MM\$)}
\label{tab:estadistica_descriptiva}
\begin{tabular*}{\textwidth}{@{\extracolsep{\fill}}
  p{1.8cm} p{1.8cm} p{1cm} p{1.8cm} p{1.8cm} p{1.8cm} p{1.5cm} p{1.8cm}}
\toprule
\textbf{Segmento} & \textbf{Variable} & \textbf{N} & \textbf{Media} & 
\textbf{Mediana} & \textbf{Desv. Est.} & \textbf{Mín.} & \textbf{Máx.} \\
\midrule
\multirow{6}{*}{DAES} 
  & P1         & 129 & 827,8     & 392,2     & 1.119,6   & 0,0    & 10.773,7  \\
  & D1         & 126 & 310,8     & 161,6     & 388,7     & 0,0    & 3.985,5   \\
  & Activos    & 117 & 4.593,1   & 2.214,5   & 5.634,2   & 0,0    & 41.370,8  \\
  & Patrimonio & 116 & 2.326,4   & 1.028,1   & 2.837,4   & 0,0    & 23.719,7  \\
  & F2         & 106 & 2.070,0   & 623,9     & 2.779,8   & 0,0    & 28.369,7  \\
  & F4         & 107 & 3.252,6   & 1.382,6   & 4.313,7   & 0,0    & 37.253,1  \\
\addlinespace
\multirow{6}{*}{CMF}  
  & P1         & 91  & 42.668,0  & 9.393,3   & 87.957,3  & 1.915,0  & 359.492,0   \\
  & D1         & 91  & 11.879,5  & 3.935,6   & 21.794,1  & 685,6    & 89.298,0    \\
  & Activos    & 91  & 359.222,7 & 54.404,1  & 787.126,6 & 12.177,0 & 3.939.576,0 \\
  & Patrimonio & 91  & 89.433,5  & 22.308,4  & 193.761,8 & 3.328,0  & 776.916,0   \\
  & F2         & 91  & 200.216,4 & 40.714,9  & 418.888,6 & 1.857,8  & 2.402.025,0 \\
  & F4         & 91  & 280.803,0 & 45.045,4  & 615.216,7 & 8.855,0  & 2.922.175,0 \\
\bottomrule
\end{tabular*}
\\[0.8em]
{\footnotesize \textit{Nota: N corresponde al total de observaciones cooperativa~$\times$~año 
con dato disponible para cada variable, sobre el período 2014--2024 para el segmento 
DAES (35 CAC) y 2013--2025 para el segmento CMF (7 CAC). Las diferencias en N entre 
variables de un mismo segmento reflejan la disponibilidad parcial de información en 
algunos estados financieros. P1~=~producción total; D1~=~remuneraciones; 
F2~=~depósitos captados; F4~=~préstamos otorgados. La alta desviación estándar 
refleja la heterogeneidad de tamaño entre cooperativas.}}
\\[0.4em]
{\footnotesize \textit{Fuente: Elaboración propia a partir de DAES y CMF.}}
\end{table}

\subsection{Requerimientos de información según los marcos internacionales}

La revisión de los tres marcos metodológicos permite identificar el conjunto 
de variables necesarias para construir una cuenta satélite del sector 
cooperativo. El detalle completo de las variables requeridas por cada marco 
---SCN 2025, Manual ONU-TSE y Manual CIRIEC--- se presenta en el 
Anexo~\ref{anx:variables-marcos}. El Cuadro~\ref{tab:sintesis} sintetiza 
qué variables son transversales a más de un marco ---y por tanto de mayor 
prioridad para la estimación--- y cuáles son exclusivas de uno solo, 
pudiendo documentarse como brechas en una primera implementación.
% ============================================================
% CUADRO 4.4 — SÍNTESIS DE VARIABLES TRANSVERSALES
% ============================================================

El cuadro siguiente cruza los tres marcos e identifica qué 
variables son exigidas por más de uno. Las variables 
transversales a los tres marcos son las de mayor prioridad 
para la estimación, ya que su ausencia comprometería 
simultáneamente la coherencia metodológica con todos los 
referentes internacionales. Las variables exclusivas de un 
solo marco tienen prioridad secundaria y pueden documentarse 
como brechas en una primera implementación.

\begin{table}[H]
\centering
\scriptsize
\setlength{\tabcolsep}{4pt}
\caption[Síntesis: variables transversales y exclusivas por marco metodológico]{%
Síntesis: variables transversales y exclusivas por marco metodológico}
\label{tab:sintesis}
\begin{tabular*}{\textwidth}{@{\extracolsep{\fill}}
  p{4.5cm} c c c c c}
\toprule
\textbf{Variable} & \textbf{Cód.} & 
  \textbf{SCN 2025} & \textbf{ONU TSE} & 
  \textbf{CIRIEC} & \textbf{Prioridad} \\
\midrule
\multicolumn{6}{l}{\textit{Transversales a los tres marcos}} \\
\addlinespace
Producción total          & P1   & \checkmark & \checkmark & \checkmark & Alta \\
Consumo intermedio        & P2   & \checkmark & \checkmark & \checkmark & Alta \\
Valor agregado bruto      & B1g  & \checkmark & \checkmark & \checkmark & Alta \\
Remuneración empleados    & D1   & \checkmark & \checkmark & \checkmark & Alta \\
Excedente bruto explot.   & B2g  & \checkmark & \checkmark & \checkmark & Alta \\
Empleo remunerado (PEP)   & PEP  & \checkmark & \checkmark & \checkmark & Alta \\
Clasificador actividad    & ---  & \checkmark & \checkmark & \checkmark & Alta \\
Cuotas de socios          & MDR  & ---        & \checkmark & \checkmark & Alta \\
\addlinespace
\multicolumn{6}{l}{\textit{Transversales a SCN 2025 y ONU TSE}} \\
\addlinespace
Transferencias corrientes & D75  & \checkmark & \checkmark & ---        & Alta \\
Transf. gobierno          & D75g\_r & \checkmark & \checkmark & ---     & Alta \\
Donaciones privadas       & D75ph\_r & \checkmark & \checkmark & ---    & Alta \\
Transf. exterior          & D75r\_r & \checkmark & \checkmark & ---     & Media \\
Renta de la propiedad     & D4   & \checkmark & \checkmark & ---        & Media \\
Formación bruta cap. fijo & P51g & \checkmark & \checkmark & ---        & Media \\
Transferencias de capital & D9r  & \checkmark & \checkmark & ---        & Media \\
Cap./nec. financiamiento  & B9   & \checkmark & \checkmark & ---        & Media \\
Depósitos captados        & F2   & \checkmark & \checkmark & ---        & Media \\
Préstamos otorgados       & F4   & \checkmark & \checkmark & ---        & Media \\
ETC empleo                & PEF  & \checkmark & \checkmark & ---        & Media \\
Horas trabajadas          & PEH  & \checkmark & \checkmark & ---        & Media \\
\addlinespace
\multicolumn{6}{l}{\textit{Exclusivas del ONU TSE}} \\
\addlinespace
Voluntarios (personas)    & VOP  & ---        & \checkmark & ---        & Media \\
Voluntarios (ETC)         & VOF  & ---        & \checkmark & ---        & Media \\
Valor econ. voluntariado  & VOV  & ---        & \checkmark & ---        & Media \\
Ingresos mercado (FEER)   & FEER & ---        & \checkmark & ---        & Media \\
Ingresos gobierno (GOVR)  & GOVR & ---        & \checkmark & ---        & Media \\
\addlinespace
\multicolumn{6}{l}{\textit{Exclusivas del CIRIEC}} \\
\addlinespace
Nombre legal              & ---  & ---        & ---        & \checkmark & Alta* \\
RUT / identificador       & ---  & ---        & ---        & \checkmark & Alta* \\
Tipo de cooperativa       & ---  & ---        & ---        & \checkmark & Alta* \\
Región o provincia        & ---  & ---        & ---        & \checkmark & Alta* \\
Número de socios          & ---  & ---        & ---        & \checkmark & Media \\
Número de beneficiarios   & ---  & ---        & ---        & \checkmark & Media \\
Tasa de morosidad         & ---  & ---        & ---        & \checkmark & Media \\
Inserción grupos vulnerables & --- & ---      & ---        & \checkmark & Baja \\
\bottomrule
\end{tabular*}
\\[0.8em]
{\footnotesize * Alta prioridad aunque exclusiva del CIRIEC, 
pues son datos de delimitación del universo sin los cuales 
no es posible construir ninguna de las cuentas anteriores.\\
\textit{Fuente: Elaboración propia en base a 
\shortcite{united_nations_system_2025}, \cite{vereinte_nationen_satellite_2018} 
y \cite{barea_tejeiro_manual_2007}.}}
\end{table}

\newpage
\subsection{Diagnóstico de disponibilidad: brechas entre requerimientos y fuentes chilenas}
\label{subsec:diagnostico-disponibilidad-ch4}  % CORREGIDO: label renombrado

\begin{table}[htbp]
\centering
\caption{Diagnóstico de disponibilidad de información por bloque}
\label{tab:diagnostico_bloques}
\small
\renewcommand{\arraystretch}{1.4}
\begin{tabularx}{\textwidth}{@{} l X X c @{}}
\toprule
\textbf{Bloque} & \textbf{Fuente chilena disponible} & \textbf{Brecha principal} \\
\midrule
1. Delimitación de la población
& Registro DAES; actividades económicas del SII; Registro Civil para vigencia legal.
& Sin variable sistemática de fecha de constitución (\texttt{FECHA\_INICIO\_VIG}: 49\,\% de faltantes); transición regulatoria de la Ley N\textsuperscript{\underline{o}}~21.641 sin registro consolidado.
 \\

2. Variables económicas del SCN
& Tramo de ventas del SII (\texttt{sales\_bracket}); estados financieros auditados para las 7 CAC supervisadas por la CMF.
& Producción solo por rangos, no valor exacto; consumo intermedio sin fuente pública (exención del Impuesto a la Renta, art.~78); remuneraciones sin desagregación.
 \\

3. Fuentes de ingreso desagregadas
& Transferencias estatales vía Ley N\textsuperscript{\underline{o}}~19.862 (0{,}57\,\% de faltantes); donaciones privadas vía registro MDS.
& Donaciones con vacíos territoriales (57{,}7\,\%) y de RUT del donante (32{,}2\,\%); sin fuente alguna para cuotas y aportes de socios.
\\

4. Trabajadores y voluntarios
& Número de trabajadores en el SII.
& Sin desagregación por jornada, sexo ni edad; trabajo voluntario inexistente en todo registro administrativo.
\\
\bottomrule
\end{tabularx}

\vspace{0.5em}
\textit{Fuente: Elaboración propia }
\end{table}

Establecido el horizonte metodológico, esta subsección contrasta los requerimientos de los manuales internacionales con la disponibilidad efectiva de información en Chile, organizando el diagnóstico en los cuatro bloques del Cuadro~\ref{tab:diagnostico_bloques}.

El diagnóstico revela un patrón consistente. En el Bloque~1, la población es identificable pero su datación es incompleta. En el Bloque~2, solo las transferencias del Estado se obtienen de forma directa: la producción se aproxima mediante el tramo de ventas del SII ---que clasifica a cada contribuyente en rangos sin revelar el importe exacto, a diferencia de España y Portugal--- y el consumo intermedio carece de fuente, ya que la exención del Impuesto sobre la Renta del artículo~78 de la Ley de Asociaciones Cooperativas elimina el registro administrativo que permitiría conocer los costos de operación.\footnote{Esta exención, si bien beneficia al sector, suprime el equivalente a la declaración del Impuesto sobre Sociedades utilizada en España para estimar el consumo intermedio. Solo las 7 CAC supervisadas por la CMF cuentan con estados financieros auditados que permiten cálculos directos.} En los Bloques~3 y~4, la brecha más estructural es la ausencia de un instrumento de encuesta directa a las organizaciones: las cuotas y aportes de socios y el trabajo voluntario ---que en Costa Rica, México, Polonia y Portugal se capturan mediante encuestas creadas para ese fin \shortcite{statistics_poland_social_2021,pedroso_conta_2023}
no tienen equivalente en Chile.

Se concluye que Chile dispone de los registros administrativos mínimos para construir una estimación del VAB de las CAC, pero que la cuenta satélite resultante opera con variables proxy en la mayoría de sus componentes. El Cuadro~\ref{tab:inventario_variables} amplía este diagnóstico con el inventario comparado de variables requeridas y su disponibilidad en los cuatro países considerados. Las brechas identificadas constituyen la agenda de datos pendiente para una versión más completa de la cuenta en el futuro \cite{archambault_lester_2022}.



\begin{table}[htbp]
\centering
\caption{Inventario comparado de variables de la cuenta satélite de CAC}
\label{tab:inventario_variables}
\small
\begin{tabular}{lllll}
\toprule
\textbf{Variable (SCN 2025)} & \textbf{España} & 
\textbf{Portugal} & \textbf{Polonia} & \textbf{Chile} \\
\midrule
Producción (P1) & 
  Declaración IS, & IES, estados & Encuesta SOF/ & 
  Estados financieros \\
  & cifra de negocios & financieros & ES-S + CIT-8 & DAES y CMF \\
\addlinespace
Consumo intermedio (P2) & 
  Declaración IS, & IES, costos & Encuesta SOF/ & 
  MIP BCCh 2018, \\
  & gastos explotación & operacionales & ES-S + CIT-8 & 
  $\alpha = 0{,}3776$ \\
\addlinespace
VAB (B1g = P1$-$P2) & 
  Derivado & Derivado & Derivado & Derivado \\
\addlinespace
Remuneración de & TGSS, & IES, costos & Encuesta SOF/ & 
  Estados financieros \\
empleados (D1) & cotizaciones & de personal & ES-S + ZUS & DAES y CMF \\
\addlinespace
Excedente de explotación & 
  Derivado & Derivado & Derivado & Derivado \\
(B2g = B1g$-$D1) & & & & \\
\bottomrule
\end{tabular}
\begin{tablenotes}
\small
\item \textit{Nota}: IS = Impuesto sobre Sociedades; TGSS = 
Tesorería General de la Seguridad Social; IES = Información 
Empresarial Simplificada (Portugal); SOF = formulario de 
organizaciones non-profit (Polonia); ES-S = encuesta de 
cooperativas (Polonia); CIT-8 = declaración tributaria corporativa 
(Polonia); ZUS = seguridad social polaco; MIP = Matriz 
Insumo-Producto; BCCh = Banco Central de Chile. 
B1g y B2g son variables derivadas en todos los países.
\item \textit{Fuentes: Elaboración propia en base a 
\cite{barea_tejeiro_manual_2007}; 
\shortcite{pedroso_conta_2023}; 
\cite{statistics_poland_social_2021}; 
\shortcite{united_nations_system_2025}.}
\end{tablenotes}
\end{table}

\newpage
\section{Métodos de estimación}
\label{sec:metodos-estimacion}

Esta sección presenta el procedimiento mediante el cual los datos administrativos disponibles se transforman en estimaciones de las variables del SCN para las cooperativas de ahorro y crédito. Se describe primero la integración de las fuentes en una base unificada, luego la operacionalización de las variables del SCN a partir de los proxies disponibles, y finalmente el modelo de cálculo del Valor Agregado Bruto y su aporte al PIB nacional, junto con los supuestos y limitaciones que condicionan la estimación.

% ============================================================
% SECCIÓN 4.3 — MÉTODOS DE ESTIMACIÓN (versión final)
% Subsecciones 4.3.1 y 4.3.2
% ============================================================

\subsection{Integración de fuentes administrativas}
\label{subsec:integracion-fuentes}

La estimación del VAB de las CAC requirió construir una base
de datos unificada a partir de tres fuentes de información
con estructuras y alcances distintos. El proceso se desarrolló
en dos etapas: primero la consolidación de cada fuente por
separado, y luego la construcción de un panel limpio con las
variables específicas necesarias para la cuenta satélite.

\paragraph{Fuente 1 — Registro institucional del universo CAC.}

A partir del registro oficial de cooperativas de la División
de Asociatividad y Cooperativas (DAES), obtenido mediante
consulta directa a su portal en mayo de 2025, se construyó
un registro del universo de análisis que abarca 38~CAC
activas del segmento no supervisado. La delimitación sigue
los criterios de elegibilidad establecidos en la
sección~\ref{sec:cuentas-satelite}: organizaciones de
propiedad democrática, con distribución limitada de
excedentes e independencia del Estado, conforme a los
criterios operacionales del Manual CIRIEC
\cite{barea_tejeiro_manual_2007}.

Este registro integra, mediante el RUT como llave de cruce,
información proveniente de cuatro fuentes adicionales: el
archivo maestro de personas jurídicas del SII, que permite
validar la existencia legal de cada cooperativa y obtener su
código de subtipo (817 para CAC); la serie de tiempo del SII,
que aporta el tramo de ventas y el número de trabajadores
para las CAC presentes en ese registro; el registro de
transferencias de la Ley N\textsuperscript{o}~19.862, que
documenta los flujos de financiamiento público hacia el
sector; y los datos institucionales de la CMF, que aportan
información sobre empleados y oficinas de las cooperativas
supervisadas. El resultado es un panel de identificación y
contexto para 35~CAC del segmento DAES y 7~CAC del segmento
CMF, sin variables económicas en el sentido del SCN.

\paragraph{Fuente 2 — Estados financieros por segmento.}

Para construir las variables económicas del SCN se
recopilaron los estados de resultados y balances generales
de cada CAC para los períodos disponibles, organizados en
dos segmentos según la procedencia de los documentos.

Para las \textbf{CAC no supervisadas por la CMF}, los estados
financieros fueron obtenidos directamente del portal DAES,
donde cada cooperativa publica sus memorias anuales. La
extracción se realizó cooperativa por cooperativa,
procesando los documentos disponibles para cada período.\footnote{Las
bases de datos consolidadas de la cuenta satélite están
disponibles en el repositorio público del proyecto:
\url{https://github.com/Iureta1/Dashboard_HuellaSocial}.} El
resultado es un panel de 133~observaciones correspondientes
a 35~CAC distintas para el período 2014--2024, con entre
2~y~24~entidades por año según la disponibilidad de memorias
en el portal.

Para las \textbf{7~CAC supervisadas por la CMF}, se
utilizaron dos fuentes complementarias: la plataforma BEST
de la CMF, que reporta datos financieros del sector
supervisado, y las memorias anuales individuales de cada
cooperativa, revisadas para obtener el desglose por entidad
que la plataforma no proporciona de forma desagregada. El
resultado es un panel de 91~observaciones que cubre las
7~CAC para el período 2013--2025, con cobertura completa en
todos los años.

\paragraph{Construcción del panel de estimación.}

A partir de las fuentes anteriores se construyó una base
limpia que contiene exclusivamente las variables necesarias
para la cuenta satélite, organizada como un panel de
observaciones CAC~$\times$~año. Cada fila corresponde a una
cooperativa en un año determinado para el que existe
información financiera disponible; las observaciones sin
datos para un período específico no se incluyen en el panel
y no se realiza imputación ni se utilizan proxies para cubrir
años faltantes. El panel consolidado del segmento DAES suma
133~observaciones y el del segmento CMF 91~observaciones,
conformando en conjunto la base de estimación de la cuenta
satélite.
% ============================================================
\newpage
\subsection{Operacionalización de variables del SCN}
\label{subsec:operacionalizacion-variables}

A partir de la base de estimación, cada variable del SCN se
construye según el procedimiento que se detalla a
continuación.

% -------------------------------------------------------
\subsubsection*{Producción total (P1)}

De acuerdo con lo establecido en la sección~\ref{subsec:cac-scn}, las CAC 
son unidades de producción de mercado cuyos ingresos operacionales constituyen 
una aproximación directa a P1 sin necesidad de valoración al costo. En 
consecuencia, P1 se obtiene directamente de la línea 
\texttt{Total\_Ingresos\_Operación} del estado de resultados de cada CAC, 
que consolida los ingresos por intereses y reajustes, los ingresos por 
inversiones y otros ingresos de operación, en línea con la definición de 
producción de mercado de una institución financiera del SCN~2025 
(párr.~7.169). El dato se extrae en pesos corrientes del año de reporte.

% -------------------------------------------------------
% -------------------------------------------------------
\subsubsection*{Consumo intermedio (P2)}

Como se estableció en la sección~\ref{subsec:cac-scn}, cuando el consumo 
intermedio no es observable en las fuentes administrativas disponibles, el 
SCN admite el uso de coeficientes técnicos derivados de matrices 
insumo-producto como procedimiento de estimación indirecta. Este es 
precisamente el caso de las CAC chilenas: la exención del impuesto de primera 
categoría establecida en el artículo~78 del DFL N\textsuperscript{o}~5 de 
2003 elimina el registro equivalente a la declaración del Impuesto sobre 
Sociedades que utilizan España y Portugal para estimar esta variable. En 
consecuencia, P2 se estima como proporción fija de la producción:

\begin{equation}
P2_{i,t} = \alpha \cdot P1_{i,t}, \qquad \alpha = 0{,}3776
\label{eq:p2}
\end{equation}

\noindent donde $\alpha$ es la razón consumo intermedio/producción del 
sector~94 (Intermediación financiera) en la Matriz Insumo-Producto de 
Chile~2018, publicada por el Banco Central de Chile. El supuesto subyacente 
es que la estructura de costos de las CAC es homogénea con el promedio del 
sector financiero. Este coeficiente se aplica de forma uniforme a ambos
segmentos ---DAES y CMF--- para garantizar la comparabilidad de los B1g
resultantes, aun cuando las CAC supervisadas publican estados financieros
auditados que permitirían una estimación directa de P2 para algunas
entidades.\footnote{El análisis de los datos contables revela que CAPUAL
y AHORROCOOP presentan razones P2/P1 implícitas de $0{,}72$, el doble del
promedio sectorial, lo que sugiere que sus gastos de administración incorporan
conceptos fuera del consumo intermedio en el sentido del SCN~2025. Estimar
P2 directamente para las entidades con desglose contable confiable
---Coopeuch, Oriencoop y Coonfía, con ratios entre $0{,}23$ y $0{,}28$---
constituye una línea prioritaria para versiones futuras de la cuenta.}
% -------------------------------------------------------
\subsubsection*{Valor agregado bruto (B1g)}
% -------------------------------------------------------

Variable derivada calculada para cada observación del panel:

\begin{equation}
B1g_{i,t} = P1_{i,t} - P2_{i,t} = (1 - \alpha)\cdot P1_{i,t}
\label{eq:b1g}
\end{equation}

\noindent El aporte agregado del sector al PIB se estima
sumando el VAB de las CAC con información disponible en cada
año:

\begin{equation}
\text{Aporte al PIB}_t =
\frac{\displaystyle\sum_{i \in \mathcal{N}_t} B1g_{i,t}}
     {\text{PIB}_t} \times 100
\label{eq:aporte}
\end{equation}

\noindent donde $\mathcal{N}_t$ es el conjunto de CAC con
P1 disponible en el año~$t$.

% -------------------------------------------------------
\subsubsection*{Remuneración de empleados (D1)}
% -------------------------------------------------------

D1 se obtiene directamente de la línea
\texttt{Remuneraciones\_y\_Gastos\_del\_Personal} del
estado de resultados. Esta línea consolida sueldos,
cotizaciones patronales y otros gastos de personal,
correspondiendo al agregado D1 del SCN~2025. En el estado
de resultados de las CAC esta cifra aparece con signo
negativo por convención contable; se utiliza su valor
absoluto. Las observaciones para las que esta línea no
está disponible en la memoria no se excluyen del panel,
pero sí quedan fuera del cálculo de B2g y del ratio
D1/B1g.

% -------------------------------------------------------
\subsubsection*{Excedente bruto de explotación (B2g)}
% -------------------------------------------------------

Variable residual calculada para las observaciones con
tanto P1 como D1 disponibles:

\begin{equation}
B2g_{i,t} = B1g_{i,t} - D1_{i,t}
\label{eq:b2g}
\end{equation}

\noindent Los impuestos sobre la producción (D2) y los
subsidios (D3) se asumen nulos. D2 es nulo en virtud de la
exención tributaria cooperativa (art.~78, DFL
N\textsuperscript{o}~5/2003). D3 se asume nulo porque los
flujos recibidos por algunas CAC bajo la Ley
N\textsuperscript{o}~19.862 ---principalmente a través del
Programa de Acceso al Microfinanciamiento--- corresponden a
convenios de prestación de servicios con el Estado, no a
subsidios a la producción en el sentido del SCN~2025.

% -------------------------------------------------------
\subsubsection*{Depreciación y VAB neto (P51d y B1n)}
% -------------------------------------------------------

La depreciación (P51d) no se incluye como variable del
cálculo central de la cuenta satélite. La razón es una
asimetría de cobertura entre segmentos: las 7~CAC
supervisadas por la CMF reportan sus gastos de apoyo en
una línea consolidada que agrupa remuneraciones,
administración y depreciaciones sin desglose, por lo que
P51d no es separable para ese segmento. Incluir B1n solo
para las CAC con memorias DAES ---donde 32~de 35~entidades
sí reportan la depreciación de forma desagregada--- generaría
una comparación inconsistente entre segmentos. En
consecuencia, el estudio reporta el \textbf{valor agregado
bruto (B1g) como variable principal}, en línea con la
práctica de Polonia en su primera iteración de cuenta
satélite cooperativa~\cite{statistics_poland_social_2021}.

A título contextual, la depreciación disponible en las
memorias DAES permite estimar que el VAB neto del segmento
DAES es en promedio inferior al bruto en un margen que se
reporta en la sección de resultados, sin incorporarse al
agregado sectorial.

% -------------------------------------------------------
\subsubsection*{Cuenta financiera: depósitos y colocaciones
(F2 y F4)}
% -------------------------------------------------------

F2 (depósitos captados) y F4 (colocaciones netas) se
obtienen directamente del balance general:

\begin{itemize}
  \item \textbf{F2} = \texttt{Total\_Depósitos}, que
  consolida depósitos a la vista y depósitos a plazo.
  \item \textbf{F4} = \texttt{Total\_Colocaciones\_Netas},
  que corresponde al total de colocaciones vigentes menos
  las provisiones sobre colocaciones.
\end{itemize}

Estas variables permiten caracterizar la escala de la
actividad de intermediación financiera de cada CAC, aunque
no forman parte del cálculo central del VAB.

% -------------------------------------------------------
\subsubsection*{Activos totales y patrimonio (Act y Pat)}
% -------------------------------------------------------

\texttt{Total\_Activos} y \texttt{Total\_Patrimonio} se
extraen directamente del balance general y se utilizan como
variables de escala para contextualizar los resultados de la
cuenta de producción.

% -------------------------------------------------------
\subsubsection*{Remanente del ejercicio (Rem)}
% -------------------------------------------------------

El remanente del período se obtiene de la línea
\texttt{Remanente\_del\_Período} del estado de resultados.
Corresponde al resultado neto después de provisiones y
corrección monetaria, y es el equivalente cooperativo del
resultado del ejercicio. Puede ser positivo o negativo.

% -------------------------------------------------------
\subsubsection*{Variables de directorio e identificación}
% -------------------------------------------------------

Las variables de directorio ---RUT, razón social, tipo de
cooperativa, región y número de socios--- provienen del
registro DAES consolidado en
\texttt{Consolidado\_cooperativas.xlsx}. Se incorporan al
panel mediante cruce por RUT. El número de socios
corresponde al dato puntual del registro DAES, sin
desagregación temporal.

% -------------------------------------------------------
\subsubsection*{Empleo remunerado (PEP)}
% -------------------------------------------------------

El número de trabajadores (PEP) se obtiene de la variable
\texttt{sii\_trabajadores} del panel DAES, que a su vez
proviene de la serie de tiempo del SII. Este dato solo está
disponible para las 12~CAC presentes en el archivo
\texttt{sii\_company\_timeseries.parquet}, que corresponde
a organizaciones con presencia en los registros de
transferencias o donaciones públicas. Para las 26~CAC
restantes no existe dato de empleo en ninguna fuente
administrativa disponible.

% -------------------------------------------------------
\subsubsection*{Variables no construibles}
% -------------------------------------------------------

Las variables que no pudieron construirse con las fuentes
disponibles quedan documentadas en el
Cuadro~\ref{tab:brechas} y constituyen la agenda de datos
prioritaria para versiones futuras de la cuenta.

\begin{table}[htbp]
\centering
\caption{Variables no construibles con fuentes actuales}
\label{tab:brechas}
\small
\renewcommand{\arraystretch}{1.5}
\begin{tabularx}{\textwidth}{@{} p{1.3cm} p{3.5cm} X @{}}
\toprule
\textbf{Cód.} & \textbf{Variable} &
\textbf{Razón de exclusión} \\
\midrule

MDR &
Cuotas de socios &
Disponible en CMF~BEST para 7~CAC supervisadas bajo la
cuenta «Capital pagado por socios» del balance. Sin
equivalente público para las 31~restantes. Variable de alta
prioridad en ONU~TSE y
CIRIEC~\cite{barea_tejeiro_manual_2007,vereinte_nationen_satellite_2018}.\\

PEF, PEH &
Empleo en equivalentes a tiempo completo y horas
trabajadas &
El SII reporta solo \textit{head count}; microdatos de
jornada no están disponibles públicamente. Brecha
estructural en todos los países de la región en primera
iteración~\shortcite{pedroso_conta_2023}. \\

D11, D12 &
Remuneración desagregada &
Las memorias anuales reportan D1 como agregado. La
desagregación entre sueldos (D11) y cotizaciones
patronales (D12) requiere microdatos tributarios de acceso
restringido. \\

VOV &
Trabajo voluntario &
Sin encuesta de uso del tiempo equivalente en Chile. Las
CAC operan con fuerza laboral remunerada; reportado como
no aplicable. \\

D75ph\_r &
Donaciones privadas &
Cruce con el registro MDS (Ley N\textsuperscript{o}~19.885):
cero registros para CAC. Exclusión por diseño legal: las
cooperativas no son elegibles por distribuir excedentes
(DFL N\textsuperscript{o}~5/2003). \\

\bottomrule
\end{tabularx}
\vspace{0.4em}
{\footnotesize
\textit{Fuente: Elaboración propia.}}
\end{table}

\subsection{Modelo de cálculo del Valor Agregado Bruto}
\label{subsec:modelo-vab}

A partir de las variables operacionalizadas en la
sección~\ref{subsec:operacionalizacion-variables}, el VAB de cada
CAC se calcula como:

\begin{equation}
VAB_{i,t} = P1_{i,t} - P2_{i,t} = (1-\alpha)\cdot P1_{i,t}
\label{eq:vab}
\end{equation}

\noindent El aporte agregado del sector al PIB se obtiene sumando
el VAB de todas las CAC identificadas:

\begin{equation}
\text{Aporte al PIB} = \frac{\sum_{i=1}^{N} VAB_{i,t}}{\text{PIB}_t}
\times 100
\label{eq:aporte-pib}
\end{equation}
\noindent La aplicación de este modelo a la base integrada produce, 
para cada CAC $i$ y año $t$, los siguientes indicadores de resultado:
(i)~producción estimada $P1_{i,t}$; (ii)~consumo intermedio imputado 
$P2_{i,t}$; (iii)~valor agregado bruto $B1g_{i,t}$; 
(iv)~remuneración de empleados estimada $D1_{i,t}$; 
(v)~excedente bruto de explotación $B2g_{i,t}$; y 
(vi)~empleo remunerado $\text{PEP}_{i,t}$. 
A nivel agregado, el resultado central es el aporte del sector 
al PIB nacional (ecuación~\ref{eq:aporte-pib}), desagregado por 
año y, cuando los datos lo permiten, por región. 
\subsection{Supuestos y limitaciones}
\label{subsec:supuestos-limitaciones}

La estimación descansa en los siguientes supuestos principales:

\begin{itemize}
    \item El tramo de ventas del SII es representativo de la
    producción real de las CAC.
    \item La razón CI/P es homogénea entre organizaciones
    del mismo tipo ($\alpha = 0{,}3776$).
    \item La remuneración media del sector financiero es
    representativa del salario promedio de las CAC.
    \item Las CAC identificadas en la DAES corresponden al
    universo activo del sector durante el período analizado.
\end{itemize}

Las principales limitaciones son la imprecisión inherente
al uso de tramos en lugar de valores exactos, la posible
subestimación de organizaciones no registradas ante el SII,
y la ausencia de datos sobre trabajo voluntario
\cite{vereinte_nationen_satellite_2018}.


\section{Panel interactivo}
\label{sec:panel}


Como complemento a los resultados presentados en los
capítulos anteriores, se desarrolló un panel de visualización
interactivo que permite explorar los datos de la cuenta
satélite de forma dinámica. El panel está disponible en línea
en el repositorio público del proyecto y fue construido en
Python utilizando la librería Plotly, generando un archivo
HTML autocontenido que no requiere instalación ni conexión
a base de datos.\footnote{\url{https://github.com/Iureta1/Dashboard_HuellaSocial}}

El panel se organiza en once pestañas temáticas que cubren
los distintos componentes del análisis. Las pestañas
\textit{Resumen CMF} y \textit{Entidades CMF} presentan la
cuenta de producción agregada y el detalle individual de las
7~CAC supervisadas para el período 2013--2025, con cifras en
MM\$ corrientes referenciadas al PIB del Banco Central. La
pestaña \textit{Segmento DAES} muestra los agregados
sectoriales anuales del segmento no supervisado, mientras
que \textit{Panel DAES} despliega el panel CAC~$\times$~año
con todas las variables de la cuenta satélite por entidad.
La pestaña \textit{Serie Histórica CMF} permite comparar la
evolución del VAB entre las 7~cooperativas supervisadas
mediante selección múltiple, y \textit{CMF vs DAES} contrasta
las variables clave entre ambos segmentos en una vista
unificada. La pestaña \textit{Distribución Regional} cruza
el universo de 35~CAC DAES y 7~CAC CMF con su localización
geográfica. Finalmente, las pestañas \textit{Cobertura},
\textit{Estadística Descriptiva} y \textit{Supuestos
Metodológicos} documentan respectivamente la proporción del
universo con datos disponibles por variable y año, las
estadísticas de resumen del panel, y los criterios
metodológicos adoptados en la estimación.

El panel constituye un instrumento de transparencia y
replicabilidad: al exponer los datos consolidados en un
formato accesible y navegable, permite verificar los
agregados presentados en la cuenta satélite, identificar
variaciones entre entidades y explorar la cobertura temporal
del panel en cada segmento.

En la trayectoria del proyecto Huella Social, este panel constituye además un avance concreto de la hoja de ruta tecnológica comprometida en el formulario FONDEF (Anexo \ref{anx:resultados tecnologicos}): materializa, para el caso piloto CAC, los reportes dinámicos y dashboards interactivos que la postulación proyecta escalar mediante una plataforma web en React con sistema API e integración de al menos 15 fuentes de datos. El panel se sitúa así como eslabón intermedio entre el prototipo exploratorio en Shiny R que antecede al proyecto (Anexo \ref{anx:A.1.1}) y la plataforma comprometida en la postulación, aportando lo que el formulario declaraba solo a nivel de diseño: una implementación operativa, documentada y verificable con datos reales del sector.